% =============================================
% Section 8: Conclusion
% =============================================
\section{Conclusion}
\label{sec:conclusion}

We have presented \tinybit{}, a framework for deploying coherent language
models on microcontrollers. Our work makes three primary contributions.

First, through systematic training and evaluation of GPT-2-style models
ranging from 1M to 15M parameters on the TinyStories dataset, we establish
an empirical lower bound on the model capacity required for coherent English
text generation. Our multi-metric evaluation framework, combining perplexity,
MAUVE scores, and LLM-as-judge assessment, provides a rigorous methodology
for characterizing the coherence-capacity trade-off under quantization.

Second, we design an Xtensa LX7-native ternary \gemv{} kernel that exploits
the \espss{}'s SIMD extensions to perform multiplication-free inference. By
encoding ternary weights as packed 2-bit representations and decomposing
matrix operations into bitwise primitives, we achieve \textbf{3.2$\times$}
speedup over na\"ive INT8 inference.

Third, we demonstrate the feasibility of \textbf{Ternary-Mamba}, a proof-of-concept
architecture combining BitNet b1.58 quantization ($\{-1, 0, +1\}$ weights) with the
Mamba state space model. At 8M parameters, Ternary-Mamba achieves a validation
loss of 2.84 (vs.\ 2.56 for FP32 GPT-2), while reducing the effective weight
precision from 32 bits to 1.58 bits---a \textbf{20$\times$ compression ratio}. This
result validates our core hypothesis that MatMul-free architectures can maintain
coherent generation at extreme quantization levels suitable for microcontroller deployment.

End-to-end evaluation demonstrates that a \textbf{1.0M}-parameter model,
quantized to INT4, generates coherent English text at
\textbf{15.82} tokens/s on the \espss{} with a total deployment footprint under
8\,MB. To the best of our knowledge, this represents the largest language
model demonstrated to produce coherent text on commodity microcontroller
hardware.

We release our training code, quantization pipeline, Xtensa kernel, and
pre-trained model weights as open-source software at
\url{https://github.com/samyeh621109/ESP-TinyLlama}, enabling the community
to reproduce our results and build upon this foundation for edge language AI.
